"\\chapter{Four banking translations before formal causal analysis}\n\\label{exec13:l2-translation}\n\nA request such as “Should we improve retention?” is not yet a causal question. The action may be a rate offer, a service contact, a product change or a budgeted policy. The unit may be an account, customer, household or branch. The outcome may be balance retention, funding cost, repayment or relationship value. An assistant should expose these alternatives and ask for the missing decision; a formal engine cannot repair a target chosen incorrectly before the graph is built. The four translations below keep the business motivation while making different estimands visible.\n\n\\section{Deposit pricing: acquisition, retention or net funding cost}\n\nSuppose a treasury manager considers a 25-basis-point deposit-rate increase. Translation A asks for the effect of offering the rate to every eligible account on thirty-day average balance. Translation B asks for the effect on balances among accounts that historically received similar offers. Translation C asks for net funding cost, including paid interest, incentive expense, repricing of existing balances and any value assigned to retained relationships. Translation D asks for a fixed-budget allocation policy that chooses which accounts receive the offer. These are four populations or outcomes, not four estimators of one effect.\n\nLet $U$ be the account, $P$ the accounts eligible under the written rule, $A$ the offer, $Y^B$ the thirty-day balance-time outcome and $Y^C$ the net funding-cost outcome. The all-eligible target is $E_P[Y^B(1)-Y^B(0)]$; the policy target additionally requires a feasible policy class and a comparator under the same budget. If a cost $C(A)$ is outside $Y^B$, expected value uses $E[Y^B(1)-Y^B(0)-C(1)]$ and subtracts that cost once. A realized post-offer balance is not automatically a pre-offer weight. A customer identifier does not establish economic ownership when accounts are shared or linked.\n\nA valid clock records when rate eligibility, customer status and balance information were available to the offer system. A statement generated after the decision cannot be used as if it were known before it. Diagnostics include overlap by eligibility stratum, repricing spillovers, concurrent campaigns and whether the offer itself changes future eligibility. A narrower overlap target can be useful, but the report must display the excluded accounts and say that the original all-eligible target remains unresolved.\n\n\\section{Arrears servicing: offer, acceptance or dynamic policy}\n\nA collections manager asks whether an early restructuring review should be offered to an account entering arrears. Translation A targets an offer at the first qualifying arrears date. Translation B targets acceptance among customers who take up the review. Translation C targets the effect of a fixed sequence of contacts. Translation D targets a dynamic policy that uses evolving hardship and payment history. Acceptance is downstream of the offer; it is not the same intervention. A static binary ATE cannot silently answer a longitudinal policy question.\n\nLet $X_t$ contain information available at the decision time, $A_t$ the offer, $L_{t+1}$ later hardship information and $Y$ discounted net recovery over 180 days. A total-effect question should not condition on a post-offer descendant $L_{t+1}$ as an ordinary baseline covariate. A policy question needs the action history, availability of each $L_t$, censoring and a comparator policy. If several accounts share a household or collection team, account rows are not automatically independent units. A field entered on day 11 cannot govern a day-10 action even if its event date is day 9.\n\nFailure cases are substantive. A model that predicts acceptance can be accurate while answering no offer-effect question. A model that trims accounts with poor overlap and continues to report an all-account effect changes the population. A missing policy-history component is an unsupported target, not a zero effect. The remedy is to return a typed alternative or refusal and request the missing policy, clock or population decision.\n\n\\section{Credit underwriting: approval threshold or feasible decision rule}\n\nA credit team may ask whether lowering an approval threshold changes loss. Translation A is the effect of offering credit to applicants between the old and new thresholds. Translation B is the effect on repayment conditional on approval. Translation C is the effect on portfolio loss including volume, utilization and capital cost. Translation D is a policy value under a capacity or risk budget. Approval is an action selected by a threshold, while repayment is observed only under selected applicants; conditioning on approved cases changes the target.\n\nThe unit is an application at the decision date. The action must specify amount, term and price, not only a binary approval label. Potential outcomes include repayment, loss, utilization and capital usage over a fixed horizon. A threshold change can alter composition, interference through capacity and future collection intensity. The information contract must distinguish bureau vintage, application-time fields, later delinquency and revised labels. Diagnostics include support around the threshold, manipulation or anticipatory behavior, policy feasibility, missing outcomes and whether the threshold is a rule or a human override.\n\nA regression discontinuity or instrumental design requires its own continuity, exclusion and assignment assumptions; a predictive score does not supply them. If the target is an approval-policy value, the report should retain the policy constraints and loss function rather than output a conditional repayment coefficient. A public or generated threshold mechanism can test translation and clock preservation before requesting bank outcome truth.\n\n\\section{Branch campaign: donor counterfactual or randomized rollout}\n\nA branch manager asks whether a campaign increased new accounts. Translation A is a single-branch time-series counterfactual using unaffected donors. Translation B is a panel effect across branches with staggered launch dates. Translation C is a cluster-randomized rollout. Translation D is the value of changing campaign intensity under a constrained budget. The outcomes, donor assumptions and uncertainty objects differ.\n\nFor a treated branch $Y_t$, donor series $X_t$ and intervention date $T_0$, a structural time-series method predicts $Y_t(0)$ from pre-period history and controls. It requires controls unaffected by the campaign and a stable relationship after $T_0$. Perfect pre-period fit does not establish either restriction. Shared marketing, customer migration or national pricing can contaminate donors. A randomized rollout addresses a different question with a different cost and design; panel methods require their own assignment and trend assumptions.\n\nAn exact construction makes the risk concrete. Let $A_t=1$ on four campaign days, $X_t=x_t^0+15A_t$ and $Y_t=100+2x_t^0+20A_t$. The pre-period relationship is $Y=100+2X$. Continuing it post-period gives a contrast of $-10$ per day, although one structural law has a true campaign effect of $20$ and a contaminated donor. A second law produces the same twelve observed rows and a true effect of $-10$ with an unaffected donor. No estimator can distinguish these laws without a restriction or experiment. With a stipulated USD1 value per application and USD20 cost, the four-day net values are +60 and -60. These are authored mechanisms, not bank evidence.\n\n\\section{The formulation record and formal handoff}\n\nThe declaration should contain $Q=(U,P,A,A_0,Y,H,\\nu,\\ell,I)$: unit, population, action, comparator, outcome, horizon, averaging measure, loss and information clock. It should preserve candidate questions when the business sentence is ambiguous. A schema check verifies fields and types; a semantic check compares them with the brief and dictionary; a formal check applies the supplied graph and method. A passed schema is not causal truth, and a failed backend is not nonidentification.\n\nThe deterministic baseline is an analyst interview plus fixed arithmetic and graph fixtures. An assisted system can retrieve evidence, draft alternatives, propose a graph and prepare typed calls. It should be scored on target fidelity, clock fidelity, preserved ambiguity, valid abstention, numerical correctness only when the target is identified, and complete failure accounting. All requests belong in the denominator, including refusals and unsupported policies. Public data and generated structural mechanisms are appropriate proposal evidence; authentic bank truth and live provider performance remain later conditional evidence.\n\nThe immediate prototype should compare the four translations against the structured baseline using fixed prompts, scenario-level splits, adverse cases and a stopping rule. A later bank study is justified only if the prototype isolates a decision-relevant uncertainty and makes its value of information plausible. No fluent report can substitute for a population decision, assignment record, clock or expert acceptance.\n"
